\section{Related Work}
\label{sec:related}

\paragraph{Orthogonalized-momentum optimizers and spectral analyses.}
\mname{Muon}~\citep{jordan2024muon} introduced polynomial Newton--Schulz orthogonalization of Nesterov momentum for 2D hidden layers while routing embeddings, norms, and the language-modeling head to AdamW.
Recent mechanistic studies have scrutinized \emph{why} orthogonalization helps: \mname{Freon}~\citep{shumaylov2026freon} showed that random or inverted singular spectra can match \mname{Muon} when global update norms and subspace alignments are controlled, motivating our explicit counterfactual norm-grafting verification (\Cref{sec:mechanism}); \citet{goyal2026spectrally} showed that amplifying small-singular-value momentum directions drives \mname{Muon}'s gains without flattening weight spectra; and \citet{dang2026batchsize} demonstrated that optimizer rankings shift with batch size and token horizon, motivating our full-budget LR sweeps and $1\text{M}$-batch and $80$-TPP scaling tests.

\paragraph{Kronecker, input-side, and two-sided preconditioning around polar maps.}
\mname{K-FAC}~\citep{martens2015kfac} factors the Fisher matrix of linear layers as $C \otimes B_{\mathrm{Fisher}}$ using model-sampled cotangents, whereas \citet{kunstner2019limitations} proved that empirical observed-label outer products $\E[\delta\delta^{\Tr}]$ reflect gradient variance rather than true Gauss--Newton curvature---a distinction we preserve throughout by treating $B_t$ empirically rather than assuming exact Fisher equivalence.
\mname{Shampoo}~\citep{gupta2018shampoo, anil2020scalable} and \mname{SOAP}~\citep{vyas2024soap} construct left and right preconditioners from full-matrix gradient Grams $\E[G G^{\Tr}]$ and $\E[G^{\Tr} G]$.
Bridging second-order preconditioning and \mname{Muon}, \mname{Newton-Muon}~\citep{du2026newtonmuon} and \mname{Toner} (\mname{MuonT};~\citealt{mohamadi2026toner}) precondition \emph{only} the input side, with \mname{MuonT} introducing decoupled $(-1/2,\,-1/8)$ inner/outer inverse roots and fast cached-eigenbasis tracking.
On the two-sided front, \mname{Mousse}~\citep{zhang2026mousse} applies symmetric $(-1/8,\,-1/8)$ gradient-Gram powers around polar orthogonalization, while \mname{GO-MUON}~\citep{che2026gomuon} applies symmetric $(-1/4,\,-1/4)$ OAS-shrunk activation/cotangent roots with global Frobenius norm grafting.
Our work shows that symmetric quarter powers with norm grafting invert at billion-token horizons and gradient Grams underperform decoupled moments, whereas combining \mname{MuonT}'s asymmetric $(-1/2,\,-1/8)$ split with width-thresholded output factorization and 1-step asynchronous execution yields consistent, scalable improvements.

\section{Conclusion and Limitations}
\label{sec:conclusion}

We investigated whether, where, and how output-cotangent second-moment preconditioning improves orthogonalized-momentum (\mname{MuonT}) optimization in Transformer pretraining.
Auditing all 11 block second moments across unpreconditioned \mname{Muon}, input-whitened \mname{MuonT}, and active \mname{WT-MLP-Lag} trajectories revealed a fundamental, optimizer-invariant architectural dichotomy: expanded $\dff$-wide pre-SwiGLU cotangents (\texttt{gate, up}) concentrate $78\%\text{--}91\%$ of their anisotropy on the diagonal while consuming $96.2\%$ of dense $O(d^3)$ eigendecomposition work, whereas $\dm$-wide cotangents (\texttt{q, k, v, o, down})---and above all residual-stream writers (\texttt{o, down}, $0.8\%\text{--}1.0\%$ diagonal anisotropy)---are rotational and drift up to $16\times$--$32\times$ faster in subspace than forward activations.
Co-designing \mname{WT-MLP-Lag} around these two laws---dense spectral roots at width $\dm$, exact per-step diagonal scaling at width $\dff$, grouped Triton backward capture with chunked row buffers, and a one-update asynchronous side-stream pipeline with non-blocking batched root fallbacks---cuts two-sided step-time overhead from $+39\%$ to $+6.8\%$ over \mname{Muon} at \texttt{130M} ($8\times\text{A100}$) and $+3.3\%$ at \texttt{520M} ($32\times\text{A100}$; $+5.6\%$ warm-state, $+6.6\%$ on $8\times\text{A100}$).
Validated against same-wave identity-factor controls (\mname{Null-WT}) and full-budget LR sweeps, \mname{WT-MLP-Lag} recovers $-0.0153$, $-0.0142$, and $-0.0093$ held-out nats against \mname{Muon} at \texttt{130M} ($3/3$ fresh paired seeds), under $80$-TPP overtraining ($10.74\text{B}$ tokens, $16/16$ Paloma domains), and zero-shot at \texttt{520M} ($15/16$ Paloma domains)---a further $28\%$, $38\%$, and $43\%$ beyond input-only \mname{MuonT}, which it also beats directly in every regime, and where pure-diagonal \mname{OD-All} falls back onto \mname{MuonT} at higher cost.

\paragraph{Limitations and future directions.}
Two limitations bound our conclusions and point to clear next steps.
First, while \texttt{130M} $1\times$ confirmation uses $3$ fresh paired seeds (and \mname{GO-MUON}/\mname{Mousse} use $2$ fresh seeds), our $10.74\text{B}$-token \texttt{130M} $4\times$ and \texttt{520M} $1\times$ runs are single-seed matched-LR trajectories due to compute budget ($\sim 25$ node-hours per \texttt{520M} run), mitigated by evaluating all $11$ second-half checkpoints and $16$ external Paloma domains.
Second, while \mname{WT-MLP-Lag} improves token efficiency by $\sim 5.3\%$, its $+6.8\%\text{--}14.7\%$ step-time overhead over \mname{Muon} across \texttt{130M} and \texttt{520M} ($8\text{--}32\times\text{A100}$; $+3.3\%\text{--}6.6\%$ over \mname{MuonT}) roughly offsets raw wall-clock savings to a fixed loss target, making its primary advantage higher final model quality at fixed token budgets or inference-amortized deployment.
As $\dm$ scales to $\ge 2{,}048$ in multi-billion-parameter models, dense $\dm \times \dm$ output roots will themselves enter the cubic bottleneck regime that $\dff$ occupied at \texttt{130M}; replacing full $\dm$ roots on \texttt{q, k, v, o, down} with head-block-diagonal ($d_{\mathrm{head}}$ or $4 d_{\mathrm{head}}$) spectral roots or sparser refresh cadences ($K=8\text{--}16$) is a natural direction for future work.
